\section{Desarrollo Experimental 7}


\section{Desarrollo Experimental 8}

En este ejercicio mediremos la respuesta en frecuencia del amplificador a lazo abierto y luego cerrado, para ello colocaremos el generador a una frecuencia de 1KHz con señal senoidal y nivel de salida anterior a la distorsión.\\
Luego de llegar a dicho punto cambiaremos de onda senoidal a cuadrada, ajustamos la base de tiempo y pendiente de disparo para obtener un solo flanco de subida, teniendo en cuenta que la punta este en x10, obtendremos la siguiente imagen para lazo abierto:

IMAGEN

Luego cerramos el lazo y obtenemos la siguiente imagen:

IMAGEN

De dichas mediciones obtenemos los siguientes datos:

\begin{equation*}
    tc (medido,lazo abierto)= AAAAAA
\end{equation*}

\begin{equation*}
    tc (medido,lazo cerrado)= BBBBBB
\end{equation*}

Para corregir este valor ya que este puede ser comparable con el tiempo de crecimiento del osciloscopio, aplicaremos las siguientes formulas:

\begin{equation*}
    tc (osciloscopio) = \frac{0.35}{AB(osciloscopio)}
\end{equation*}

y

\begin{equation*}
    tc (amplificador) = \sqrt{tc(medido)^2 - tc(osciloscopio)^2}
\end{equation*}

Otenemos entonces los siguientes en ancho de banda:

\begin{equation*}
    BW (lazo abierto) = \frac{0.35}{tc(lazo abierto)}
\end{equation*}
\begin{equation*}
    BW (lazo cerrado) = \frac{0.35}{tc(lazo cerrado)}
\end{equation*}